\section*{Abstract}
Zero-shot cross-lingual named entity recognition (NER) depends heavily on the choice of source language; however, it is not settled what makes a good source for South African Bantu languages. This study fine-tunes AfroXLMR-base on eleven sources from two corpora, NCHLT (government text) and MasakhaNER~2.0 (news), and evaluates each model on nine test sets in six languages, which gives 297 zero-shot and 36 few-shot evaluations over three seeds. Linguistic relatedness is the most consistent predictor of transfer ($+0.10$ transfer ratio per relatedness step, $p<.001$), and it remains significant when corpus terms and the tokenizer-fertility gap are added. However, corpus match can override relatedness: a sister-language news source beats a same-language government source by 32.0~F1 on isiZulu news and by 12.7~F1 on isiXhosa news (Holm $p=.005$), but not on isiXhosa or Setswana government text. Targets seen in AfroXLMR's pretraining do not transfer better than unseen targets, although exposure is confounded with language cluster. With 500 target sentences, 57--89\% of the gap to in-language performance closes. Contrary to our hypothesis, Swahili is a stronger Setswana source than Sepedi or Sesotho by 13.0--19.0~F1, and it stays ahead when every source is trained on the same number of entity spans. Since the two corpora differ in both text domain and annotation guidelines, we close by proposing experiments that could separate the two.

\smallskip
\noindent\textbf{Keywords:} cross-lingual transfer; named entity recognition; South African languages; low-resource NLP; AfroXLMR
